\begin{document}
\maketitle

\begin{abstract}
Quantization changes sparse autoencoder (SAE) activations, but an activation change need not change what an intervention on that feature does. In Gemma-2-2B with a frozen layer-12 Gemma Scope SAE, 27 FLORES-200 languages and round-to-nearest weight quantization, we first show that activation survival, language-association retention and causal contribution are different outcomes: 70.8\% of established language features keep their peak language at INT4, removing 64-feature language groups raises own-text loss beyond matched controls in all 27 languages, and individual feature effects drift as much for matched other directions. We then cross the two things quantization changes, the activations before an intervention (context) and the weights after it, for a fixed feature ablation. At INT6 the effect is preserved to within about two numerical floors (cosine 0.98). At INT4 the group-level harm is amplified (median net harm 0.85 against 0.26), mostly through the context ($|C|=0.31$ against $|W|=0.10$; context larger in 19 of 27 languages), while the effect vector of single features is rotated (cosine 0.77), mostly by the weights, with a substantial interaction and not specifically for the feature's own direction. Averaged fingerprints overstate transfer (+0.126 cosine). The sensitivity to other directions predicts the held-out error (21.5\% reduction at INT4, but not at INT6, and equally in a rounding-only state); activation survival and gradient attribution add nothing. Apart from one family, all analyses are exploratory; one model, one SAE, simulated quantization.
\end{abstract}

\section{Introduction}
Weight quantization reduces the numerical precision of language-model parameters, and its effects need not be uniform across languages or tasks \citep{marchisio-etal-2024-quantization}. Internal measurements can complement task metrics, but only if their interpretation is explicit. Calling a sparse autoencoder (SAE) feature ``lost'' can mean a changed activation pattern, a changed preferred language, or a changed effect on predictions. SAEs decompose an activation into sparse directions \citep{cunningham2023}; language-associated features have been identified and manipulated in full precision \citep{deng-etal-2025-unveiling}; and \citet{duan2026} measure activation correlations before and after weight quantization with a frozen SAE, leaving open whether interventions built on those features survive. The representation can change while a contribution remains, or an intervention can change because the downstream model changes.

We began by asking whether language-selective features were more vulnerable than non-selective controls and whether that explained resource-tier disparities. After implementation audits and negative primary contrasts we expanded the project to test the measurement assumptions directly, and then to its main question. This report presents the whole investigation, including hypotheses that were not supported. \textbf{RQ1:} how do activation preservation and language-association retention compare? \textbf{RQ2:} do selected language groups contribute to prediction and comprehension, and how does that change with precision? \textbf{RQ3:} do individual intervention effects support an interpretation of feature-specific damage? \textbf{RQ4:} does the group dependence appear on a second model? \textbf{RQ5 (main):} does the effect of a \emph{fixed} intervention survive quantization, and is a change due to the context, the weights, or their interaction; can it be predicted for held-out features?

Our contribution is a controlled case study: native interventions, donor patches and effect fingerprints across weight precisions (Sections 5 and 6), and a crossed context--weights design for a fixed feature intervention with matched and random directions, a numerical floor, a held-out prediction target and an external-text check (Section 7). Ablation, patching and fingerprints are not new methods; the outcome is a set of distinctions and counterexamples to reading activation survival as functional survival, within one model and one measurement basis.

\section{Related Work}
\paragraph{Multilingual quantization.}
\citet{marchisio-etal-2024-quantization} show that task, language and evaluation method affect the observed degradation; their behavioural evidence motivates our language panel but does not identify SAE features as its mechanism. Resource categories are a descriptive grouping \citep{joshi-etal-2020-state}, not controlled interventions on training data.

\paragraph{Language-associated features.}
\citet{deng-etal-2025-unveiling} identify language-specific SAE features and use ablation and steering, so an own-language loss effect in full precision is supporting evidence, not a novelty claim. Gemma Scope supplies open SAEs across Gemma 2 layers \citep{lieberum-etal-2024-gemma}. A frozen dictionary gives consistent coordinates across precision states without guaranteeing that each coordinate is an invariant mechanism.

\paragraph{Compression and causal interpretation.}
\citet{duan2026} measure feature survival under round-to-nearest quantization by activation correlation. A pruning study already separates geometric survival from causal importance \citep{borobia2026}. \citet{chang-etal-2025-inputs} use activation patching to localize quantization error and split it between inputs and weights at the loss level, which is why a context-dominated result alone is not our contribution; \citet{hoang2026} analyze the geometry of intervention-induced logit changes; and \citet{vaidyanathan2026} show that the effect estimated by activation patching contains an interaction term that is a diagnostic rather than noise, small when the model is locally affine. We add a crossed design for a fixed SAE intervention under quantization, controls for generic drift, and a held-out prediction target. Belebele \citep{bandarkar-etal-2024-belebele} supplies a separate comprehension outcome.

\section{Setting and Research History}
We use Gemma-2-2B \citep{gemma2024} and a 16,384-feature JumpReLU \citep{rajamanoharan2024} Gemma Scope SAE at the layer-12 post-block residual stream; the model runs in bfloat16, the SAE and all edits in float32, and the same dictionary serves every condition. Quantization uses the vendored reference round-to-nearest (RTN) functions (output-channel scaling), restoring the original weights before each state; values stay in floating-point storage, so we do not evaluate integer kernels. States are full precision (FP), INT8 to INT4, and INT16 and INT12 as near-lossless rounding references (re-rounding bfloat16 weights still changes them). One BOS token is prepended; BOS and padding are excluded from statistics and edits.

The panel is 27 parallel FLORES-200 \citep{nllb2022} languages in three descriptive tiers of nine (Appendix~\ref{app:languages}). Features are discovered on the development split; association studies use 27,036 content tokens per language on devtest; the loss evaluation uses 300 aligned devtest sentences per language. Devtest had been examined earlier, so all FLORES functional analyses are exploratory. Belebele items are split by passage origin into calibration and evaluation; calibration selects the answer format and 18 languages with accuracy above chance plus 0.05, and the frozen evaluation has 385 items per language.

An initial implementation was retired after an audit; its outputs are not used. A rule is \emph{pre-specified} if committed to the repository before the data it judges were scored; none was registered externally. The only confirmatory family is the Belebele family (its first statement needs a positive FP macro effect; 99\% intervals for the five retention statements); everything else is exploratory.

\section{Measurements and Interventions}
\paragraph{Activation and association.}
With $a_f^b(t)$ the activation of feature $f$ at token $t$ in state $b$, the preservation score is $c_f^b=\operatorname{corr}_t(a_f^{\mathrm{FP}}(t),a_f^b(t))$ and the historical survival cutoff is $c_f^b>0.9$. Language selectivity uses firing rates $r_{\ell f}$ and $S_f=\max_\ell r_{\ell f}/\operatorname{mean}_\ell r_{\ell f}$, with a shuffled-language null (95th percentile, 50 permutations), at least 50 peak tokens, a panel rate of at least $10^{-5}$ and the same peak in 80 of 100 resamples; the cohort has 8,086 features. A feature is \emph{established} on devtest if its original language is the unique supported firing-rate peak, and \emph{retained} at $b$ if that stays so; retention conditions on FP establishment. We also analyze the FP margin to the strongest competing language.

\paragraph{Groups, controls and edits.}
$G_\ell$ is the 64 highest-activation selective features of language $\ell$ on dev. Controls are matched greedily on log firing rate and log mean activation from 2,860 supported non-selective features (five draws; three in Study G). Matching is imperfect (75 of 81 language--size ratios within $[0.8,1.25]$; Thai, Yoruba and Lao outside at K64), so activity matching is not an identification argument for language specificity. With $x$ a residual, $a$ its SAE encoding and $d_f$ a decoder row, edits preserve the reconstruction error:
\begin{equation}
 x'=x+\sum_{f\in G}(a'_f-a_f)d_f,
 \label{eq:edit}
\end{equation}
zero ablation sets $a'_f=0$ and scaling $a'_f=\alpha a_f$; identity and sham edits check the path. For donor patches the recipient is quantized and $a'_f$ is the FP activation at the same position.

\paragraph{Effects and uncertainty.}
Harm is the change of the negative log-likelihood against the same state's unedited prediction (nats per token; Belebele nats per item). The group effect is
\begin{equation}
 E_\ell^b=h_\ell^b(G_\ell)-\frac{1}{R}\sum_{r=1}^{R}h_\ell^b(C_{\ell r}),
 \label{eq:effect}
\end{equation}
negative harm is improvement and never clipped. Retention is $E_\ell^b/E_\ell^{\mathrm{FP}}$ for languages with a clearly positive FP effect, averaged per language (a mean of ratios; it differs from the ratio of means). Loss intervals are 2,000 paired cluster-bootstrap draws over aligned sentence identifiers; Belebele resamples passages; association intervals use 1,000 sentence resamples and can be shifted by near-tied peaks.

\section{Replication and Initial Tests}
The vendored WikiText replication matches the published table closely (INT7 81.19\% against 81.26\%; INT6 51.15\% against 51.30\%), but the no-BOS FP perplexity is 458.4; with BOS added and masked it is about 16.1 and INT6 survival is near 97.5\%. This shows sensitivity to the input regime, not that every published result is invalid, so we use BOS-masked evaluation and keep the no-BOS study as history. In that primary regime none of the three planned tests rejects after Holm correction (selective-minus-control survival $-0.28$ and $-1.44$ points at INT7 and INT6; adjusted $p=0.514$, $0.128$); a stricter post hoc threshold gives $-7.16$ at INT6, and matching activation magnitude reverses it to $+4.58$. The tier disparity (INT6 survival 69.4/90.6/95.1\% for low/mid/high) attenuates from $+12.3$ to $-0.81$ points per step after conditioning on activation scale ($p=0.667$), which is correlational and does not identify the cause of multilingual damage.

\section{Association and Functional Results}
\paragraph{Peaks persist beyond correlation survival.}
Only 6,663 of 8,086 selected features (82.4\%) establish their peak on FP devtest; among them, retention falls gradually to 70.8\% at INT4 (Table~\ref{tab:ladder}); the INT16 reference loses 56 features (99.2\% retained) and no established feature becomes inactive. Correlation and association disagree: at INT6, 75.8\% of association-lost features still exceed the correlation cutoff; at INT4, 79\% of retained ones do not. The FP margin predicts association loss within a language (AUROC 0.90, 0.81, 0.75 at INT6/5/4; most borderline fifth against least at INT6: 20.7\% against 0.1\%), so much of the loss is sensitivity to near ties.

\begin{table}[t]
\centering
\begin{tabular}{lrrr}
\toprule
State & Corr. (\%) & Assoc. (\%) & Effect ratio\\
\midrule
INT8 & 100.00 & 98.1 & 1.01\\
INT7 & 99.97 & 97.5 & 1.03\\
INT6 & 94.94 & 94.7 & 1.04\\
INT5 & 59.00 & 89.1 & 1.13\\
INT4 & 13.75 & 70.8 & 2.14\\
\bottomrule
\end{tabular}
\caption{Different measurements across precision. Correlation survival pools active features; association retention conditions on 6,663 established features; the effect ratio is the mean of 27 per-language ratios. Denominators and meanings differ; all exploratory.}
\label{tab:ladder}
\end{table}

\paragraph{Group dependence on loss, and little localized damage.}
Zeroing a language's K64 group raises own-text loss beyond controls in all 27 languages at every precision: FP macro effect 0.704 nats/token $[0.684,0.723]$ (median 0.262), retention near one through INT6, 1.13 at INT5 and 2.14 at INT4 $[2.04,2.50]$ (the ratio of macro effects is only 1.33). This is dependence of a damaged model on the group, not better language ability. Restoring the group with FP activations at INT4 recovers 2.6\% of the loss gap (controls 1.2\%; all SAE features 17.4\%; whole residual 39.3\%), and about nothing at INT5 and INT6: necessity within a model is not responsibility for the difference between models. Matching association-lost and retained features on native activity shows no systematic harm difference (median lost-minus-retained $+0.0006$, $+0.0007$, $+0.0002$ nats/token at INT4/5/6).

\paragraph{The comprehension statement is not met.}
On frozen Belebele the FP macro group effect is $-0.0001$ nats/item $[-0.0035,+0.0031]$; accuracy is 0.426 unedited, 0.426 with the group ablated and 0.425 with controls. No covered language has a clearly positive FP effect, so statement S0 is not met and S1--S5 are undefined; the whole group ($+0.0007$ $[-0.0030,+0.0042]$) agrees. Random sets of 1,024--2,860 features lower accuracy by 1.2--6.0 points while sham edits change nothing, so the edit path is sensitive and the null cannot exclude group effects of about half a point. Our frozen prompt scores above the lm-evaluation-harness prompt in all 18 languages (point estimates).

\paragraph{Individual effects drift, but directions drift too.}
For 1,487 features the median cosine between FP and quantized effect fingerprints is 0.997/0.988/0.961/0.839 at INT8/6/5/4 (fraction above the INT16 reference 94.0/89.7/66.5/0.3\%), yet each feature's effect stays far closer to its own FP fingerprint than to others' (retrieval 0.997--1.000). In the 404-feature control that subtracts another matched decoder direction at the same positions and amplitudes the medians are 0.842 and 0.828 at INT4 (differences of medians 0.004--0.007 at INT8/6/5), so most drift is shared by arbitrary perturbations. About 31\% of features above the write-side reference fail the activation-survival cutoff at INT5, and association status has descriptive AUROC 0.506 for INT4 effect similarity (the FP margin predicts association loss at 0.86 but function change at 0.58): label stability and causal-effect stability are different.

\begin{figure}[t]
\centering
\includegraphics[width=\linewidth]{direction-control.pdf}
\caption{Median fingerprint similarity on the same 404 features and matched directions. Bars show interquartile ranges, not confidence intervals. The numerical drift is largely shared by alternative perturbations.}
\label{fig:directions}
\end{figure}

\section{Intervention Transfer under Quantization}
\label{sec:transfer}
\paragraph{Design.}
The comparisons above change activations and weights together. We instead ask how the effect of a \emph{fixed} intervention changes when the model is quantized, and whether the change comes from the context (the residual and SAE activations before the edit), the weights (the layers after it) or their interaction. At the positions where a feature is active in FP we subtract its FP decoder contribution (float32; every cell subtracts its own unedited baseline; a zero edit is exactly zero). Let $D[p,s]$ be the mean centered change of the output logits with prefix state $p$ and suffix state $s$, each FP or quantized (Q). Then $C=D[\mathrm{Q},\mathrm{FP}]-D[\mathrm{FP},\mathrm{FP}]$, $W=D[\mathrm{FP},\mathrm{Q}]-D[\mathrm{FP},\mathrm{FP}]$, $I=D[\mathrm{Q},\mathrm{Q}]-D[\mathrm{Q},\mathrm{FP}]-D[\mathrm{FP},\mathrm{Q}]+D[\mathrm{FP},\mathrm{FP}]$, and $D[\mathrm{Q},\mathrm{Q}]-D[\mathrm{FP},\mathrm{FP}]=C+W+I$ exactly. A fifth cell uses the quantized state's own activation (native amplitude).

\paragraph{Controls, floor, populations.}
Each edit is repeated along a matched swapped direction (another feature's decoder row) and three random norm-matched orthogonal directions. Because the model runs in bfloat16, even INT16 changes the logits; per direction kind the numerical floor is the 95th percentile of the INT16 error over calibration features (28.7 for the real direction), and an item is \emph{eligible} if its FP effect norm exceeds five floors. The normalized error is $\lVert D[\mathrm{Q},\mathrm{Q}]-D[\mathrm{FP},\mathrm{FP}]\rVert/(\lVert D[\mathrm{FP},\mathrm{FP}]\rVert+\text{floor})$. The feature level has 404 features with up to 16 contexts each (6,464 feature--context pairs; 2,203 eligible real-direction items of 235 features), split by a hash of the feature id into calibration and held-out halves; the group level ablates each language's K64 group and three matched control groups on the 300 evaluation sentences of the 27 languages and decomposes the net harm of Equation~\ref{eq:effect}. Intervals are 2,000-replicate crossed bootstraps over features and source-sentence clusters (group level: sentences only). The design, thresholds and hypotheses were frozen and hashed before any quantized outcome was read; the claim rules were tightened once after an independent code review and before the outcomes were read (Appendix~\ref{app:rules}).

\begin{table*}[t]
\centering\small
\begin{tabular}{lrrrrrrrrrr}
\toprule
 & \multicolumn{4}{c}{Group level (27 languages, nats/token)} & \multicolumn{6}{c}{Feature level (2,203 eligible items)}\\
\cmidrule(lr){2-5}\cmidrule(lr){6-11}
State & net harm & $|C|$ & $|W|$ & $|I|$ & cosine & error & floors & $C$ & $W$ & $I$\\
\midrule
INT16 (rounding) & 0.262 & 0.0005 & 0.0003 & 0.0002 & 0.998 & 0.056 & 0.71 & 0.055 & 0.053 & 0.074\\
INT6 & 0.281 & 0.009 & 0.009 & 0.001 & 0.984 & 0.165 & 1.9 & 0.100 & 0.145 & 0.086\\
INT4 & 0.848 & 0.310 & 0.104 & 0.071 & 0.768 & 0.664 & 7.7 & 0.359 & 0.560 & 0.315\\
\bottomrule
\end{tabular}
\caption{Main results of the transfer study. Group level: median net harm with the FP edit (FP: 0.263) and median absolute terms over 27 languages. Feature level: medians over eligible real-direction items of the cosine between the quantized and FP effect vectors, the normalized error, the error in numerical floors, and the norms of $C$, $W$, $I$ relative to the FP norm plus floor. INT16 is the rounding-only reference.}
\label{tab:transfer}
\end{table*}

\begin{figure*}[t]
\centering
\includegraphics[width=\textwidth]{transfer-terms.pdf}
\caption{(a) Median relative norm of the context, weight and interaction terms of the effect vector, real direction (with 95\% interval) and the four control directions. (b) Group-level absolute context and weight terms of the net harm, one point per language. (c) Cosine of the averaged fingerprint minus the mean cosine of its contexts, with the rounding-only state as reference; the interval resamples features only.}
\label{fig:terms}
\end{figure*}

\paragraph{Group level.}
At INT6 the terms are small and equal (median $|C|=0.0091$ $[0.0056,0.0114]$, $|W|=0.0092$ $[0.0061,0.0133]$; the context term is larger in 12 of 27 languages), 3\% of the FP net harm and 10 to 30 times the rounding-only values, so the group effect is nearly preserved (net harm 0.281 against 0.263). At INT4 the median net harm is 0.85 with the FP edit and 0.43 with the quantized state's own activation, and the context term is about three times the weight term ($|C|=0.310$ $[0.258,0.334]$, $|W|=0.104$ $[0.094,0.122]$, $|I|=0.071$); the interval of the median $|C|-|W|$ is $[+0.022,+0.073]$ and the context term is larger in 19 languages $[16,21]$. The expectation that context dominates (A3, suggested by loss-level splits) is supported at INT4 and undecided at INT6.

\paragraph{Feature level.}
At INT6 the effect vector is rotated little (cosine 0.984, error 16.5\% of the effect, 1.9 floors); at INT4 it is rotated strongly (cosine 0.768) with a slightly larger norm (ratio 1.09) and an error of 66\% $[0.647,0.680]$, 7.7 floors. The error is almost orthogonal to the FP effect (97--99\% at every state, 99.7\% at INT16) and using the quantized state's own activation leaves 97\% of it, so neither a rescaling nor an activation change explains it (H1: the shares alone say nothing because they are equal in the rounding-only state, hence the requirement that the error clears the floor). Unlike the group-level loss, the \emph{weight} term exceeds the context term in the effect vector (relative norms 0.145 against 0.100 at INT6, 0.560 against 0.359 at INT4; $\lVert C\rVert>\lVert W\rVert$ in 6\% and 10\% of items), and the interaction is a third of the change at INT4. The two readings differ in what they measure: the norm of a mean logit change against a signed loss summed over a group. The error is not specific to the feature: matched and random directions give similar terms and errors (INT4 gap of the real direction over its controls $+0.019$ $[0.009,0.032]$, negative at INT6; the controls are not eligibility-filtered, so the contrast mixes scale and specificity), in line with the generic drift of Section 6.

\paragraph{Averages hide the contexts.}
The cosine of a feature's \emph{averaged} fingerprint exceeds the mean cosine of its single contexts by 0.126 $[0.121,0.131]$ at INT4 (404 of 404 features), 0.024 at INT6 and 0.0145 $[0.011,0.017]$ in the rounding-only state, where independent noise averages out (Figure~\ref{fig:terms}c); the contexts of a feature agree in FP only with coherence 0.59. Averaged fingerprints therefore overstate how well an intervention transfers (H4).

\begin{figure}[t]
\centering
\includegraphics[width=\linewidth]{transfer-predict.pdf}
\caption{Relative reduction of the held-out error of the model with the generic sensitivity (A+B) against the strongest baseline, with crossed-bootstrap intervals; dashed: the 10\% criterion.}
\label{fig:predict}
\end{figure}

\paragraph{Prediction.}
We fit ridge models on the calibration half and score the held-out features (feature-held-out: their contexts share source sentences with the calibration half). A uses activation measurements (cosine of FP and quantized activations, log activation, firing share, log effect and residual norms), G adds the loss gradient at the edited rows, B adds the generic sensitivity measured with the \emph{other} four directions at the same context. B reduces the held-out error against the strongest baseline (A+G) by 21.5\% $[12.7,29.4]$ at INT4, so the 10\% criterion is met, but by 10.9\% $[-1.2,21.4]$ at INT6, where the baselines are not better than the calibration median; the pre-specified utility criterion needs both states and is not met (Figure~\ref{fig:predict}). The same model reduces the error by 23.4\% $[15.1,30.8]$ in the rounding-only INT16 state, so B measures how noise-prone a context is rather than a quantization-specific mechanism. Neither gradient attribution (INT4 error 0.124 against 0.125) nor the Pearson survival of the feature adds to A or to B (INT4: $-0.4\%$ $[-1.8,0.7]$ and $-0.6\%$ $[-4.3,3.4]$).

\paragraph{External panel.}
\textbf{[PENDING: six-language Wikipedia panel, group and feature level, same rules.]}

\section{Graded Feature Scaling}
Study G varies $\alpha$ in Equation~\ref{eq:edit} for the K64 groups with three down multipliers (0.3536, 0.1768, 0.0884) and their reciprocals, calibrated on a 100-sentence FP pilot, on 300 sentences, 27 languages and three control draws. In FP the median net harm rises 0.0223, 0.0902, 0.1687, 0.2624 as the group is suppressed more strongly and then zeroed (ordered in all 27 languages; 26 with a positive lower bound at the deepest down multiplier); strong amplification also raises harm; the curves are similar from FP through INT5 with an INT4 increase. The strict precision-monotonicity rule is not met and no relaxed tolerance is substituted. A follow-up on 562 reserve sentences with five control draws finds a small INT4 benefit of amplifying by 1.1 (median net harm $-0.00135$ nats/token, 97.5\% interval $[-0.00197,-0.00032]$, 20 of 27 languages; 1.2: not met), with a small harm at FP, INT6 and INT5. State-defined groups (Jaccard with the FP group 0.49 at INT4) are not more consequential ($-0.0098$), and the language-specific optimal gain is reproducible (reliability 0.917) but not explained by inverse activation shrinkage (rank correlation $-0.170$, $p=0.804$): a simple compensation explanation is not supported.

\begin{figure}[t]
\centering
\includegraphics[width=\linewidth]{gain-down.pdf}
\caption{Median net harm under suppression, from the completed Study G curves. Net harm subtracts controls scaled by the same multiplier. The strongest precision change occurs at INT4.}
\label{fig:gain}
\end{figure}

\section{A Second-Model Check}
\label{sec:llama}
We ran a minimum consistency check on Llama-3.2-1B \citep{grattafiori2024} with a public TopK \citep{gao2024} SAE (14,336 features, residual entering block 12) through an adapter around the frozen Gemma pipeline (same languages, sentences, selectivity rule, K64 groups, three control draws, Study G multipliers, 300 sentences; FP, INT8, INT6, INT5, INT4), with predictions fixed beforehand (Table~\ref{tab:llama}). The verdict is \emph{not consistent}: the group effect exists but is small (median $0.048$ nats/token against $0.262$; positive in 21 of 27 languages, 24 required; large in German, Yoruba and Amharic, negative in six languages); it is flat through INT6 (P2a met) but without INT4 amplification (the frozen mean retention is $-3.57$ because one language with $E=0.0027$ has a ratio of $-88.9$; the post hoc median is $0.86$, against $1.54$ on Gemma), and the down-dose ordering holds in 20 of 27 languages at FP and 15 at INT4. Control draws are heavy-tailed, so post hoc robustness variants cannot change the verdict. Model family, SAE architecture, depth and tokenizer all differ, and this dictionary labels 91\% of its supported features selective with only 898 candidate controls, so the check does not show that the Gemma findings fail elsewhere.

\begin{table*}[t]
\centering\small
\begin{tabular}{p{4.2cm}p{4.9cm}p{3.5cm}p{1.1cm}}
\toprule
Prediction (fixed beforehand) & Llama-3.2-1B & Gemma-2-2B & Met\\
\midrule
P1: median $E>0$ with its interval above 0 and $E>0$ in $\geq 24$ of 27 languages & median $0.048$ $[0.018,0.173]$; 21 of 27 & median $0.262$; 27 of 27 & no\\
P2a: macro retention in $[0.85,1.15]$ at INT8 and INT6 & $1.010$, $1.024$ & $1.006$, $1.038$ & yes\\
P2b: INT4 retention $\geq$ INT6 $+0.30$ and the largest & INT5 $0.14$, INT4 $-3.57$ (median of ratios $0.86$) & $1.13$, $2.14$ (median $1.54$) & no\\
P3a: down-dose ordering (H-G1) in $\geq 24$ of 27 languages and in the median, at FP and INT4 & 20 of 27 at FP, 15 of 27 at INT4 & met (Study~G) & no\\
P3b (reported only): the deepest down multiplier costs more at INT4 in $\geq 20$ of 27 and the median ratio exceeds 1.3 & 14 of 27; median ratio $1.17$ & ratio $1.91$ & no\\
\midrule
Verdict & \multicolumn{3}{l}{not consistent (P1 fails)}\\
\bottomrule
\end{tabular}
\caption{Pre-specified predictions and outcomes for the Llama-3.2-1B check. $E$ is the group effect of Equation~\ref{eq:effect} (nats/token); retention is the mean of per-language ratios $E^b_\ell/E^{\mathrm{FP}}_\ell$ over languages with a clearly positive FP effect (21 of 27). The median of the ratios is post hoc and not used for the verdict.}
\label{tab:llama}
\end{table*}

\section{Discussion and Conclusion}
The original selective-destruction account is not established: primary contrasts were weak, the threshold-dependent disadvantage reversed with magnitude matching, and donor patches attributed little quantization loss to the selected groups. The main methodological lesson is to say what a survival measurement preserves: a correlation cutoff describes activation agreement, a peak label depends on a margin, and a logit fingerprint describes a sampled intervention in a changing downstream model.

The crossed design adds three findings. First, whether an intervention keeps its effect depends on the precision more than on the feature: it does at INT6 within about two floors and does not at INT4, where the effect vector is rotated (cosine 0.77) rather than shrunk. Second, the answer to ``context or weights'' depends on the quantity: the group-level loss is amplified mostly through the context, the effect vector is rotated mostly by the weights, and the interaction is never negligible at INT4, so a single split would mislead. Third, the transfer error is a generic response of the model, shared by matched and random directions and predicted by how noise-prone a context is, not by activation survival or by gradient attribution; and averaged fingerprints hide it. Together with the earlier result that activation survival and association status do not track functional change, this argues against using activation survival as a proxy for the survival of interventions; it does not show that a better proxy is unavailable.

\label{main-content-end}
\section*{Limitations, Challenges, and Future Work}
All analyses use one base model, layer, frozen SAE and simulated round-to-nearest weights (not deployed integer kernels); dictionary fidelity can deteriorate after compression. Zero ablation can be off-distribution and does not establish unique necessity or semantic identity; matching is imperfect and control draws overlap. Language tiers confound script, tokenization, training data and competence. Most results are exploratory, and bootstrap intervals condition on the selected sample.

For the transfer study: the numerical floor is large (bfloat16 execution), so INT6 effects are only about two floors; ``held-out'' means held-out features (88\% of held-out source sentences also occur among calibration contexts); the group-level interval resamples sentences only and the averaged-fingerprint interval features only; the control directions of the specificity contrast are not eligibility-filtered; the ridge strength chosen was the largest of the three tried; the generic-sensitivity model also works in the rounding-only state, so its INT4 gain must not be read as quantization-specific; the gradient baseline is the loss-attribution form, not the effect-vector form; and the six-language external panel is a validity check, not a replication of the 27-language study. The second-model check confounds model family, dictionary, depth and tokenizer; Gemma-2-9B with Gemma Scope was handed to another team and is not part of this report.

Challenges included replacing the initial pipeline, accelerator memory limits, hardware heterogeneity (one node whose float computation differed and one that ran the model on the CPU, both caught by gates and log checks), and CUDA failures. Future work should replicate across models and quantizers (including integer kernels and activation quantization), test behavioural steering with calibrated headroom, and test the generic-sensitivity predictor against a quantization-specific alternative. These are not completed results.
